\item Sauve-moi, m\frflex{o}n Dieu :\psstar{} les eaux m\frflex{o}ntent jusqu’à ma gorge !
\item J’enfonce dans la v\frflex{a}se du gouffre,\psstar{} ri\frflex{e}n qui me retienne ; 
\item Je descends dans l’ab\frflex{î}me des eaux,\psstar{} le fl\frflex{o}t m’engloutit.
\item Je m’épuise à crier, ma g\frflex{o}rge brûle.\psstar{} Mes yeux se sont usés d’att\frflex{e}ndre mon Dieu.
\item Plus abondants que les chev\frflex{e}ux de ma tête,\psstar{} ceux qui m’en ve\frflex{u}lent sans raison ;
\item Ils sont nombre\frflex{u}x, mes détracteurs,\psstar{} à me ha\frflex{ï}r injustement. 
\item Moi qui n’ai rien volé, que devr\frflex{a}i-je rendre ?\psstar{} Dieu, tu connais ma folie, mes fautes sont à n\frflex{u} devant toi.
\item Qu’ils n’aient pas honte pour m\frflex{o}i, ceux qui t’espèrent,\psstar{} Seigneur, Die\frflex{u} de l’univers ;
\item Qu’ils ne rougissent pas de m\frflex{o}i, ceux qui te cherchent,\psstar{} Die\frflex{u} d’Israël !
\item C’est pour toi que j’end\frflex{u}re l’insulte,\psstar{} que la honte me co\frflex{u}vre le visage :
\item Je suis un étrang\frflex{e}r pour mes frères,\psstar{} un inconnu pour les f\frflex{i}ls de ma mère.
\item L’amour de ta mais\frflex{o}n m’a perdu ;\psstar{} on t’insulte, et l’insulte ret\frflex{o}mbe sur moi.
\item Si je pleure et m’imp\frflex{o}se un jeûne,\psstar{} je reç\frflex{o}is des insultes ;
\item Si je revêts un hab\frflex{i}t de pénitence,\psstar{} je deviens la f\frflex{a}ble des gens :
\item On parle de m\frflex{o}i sur les places,\psstar{} les buveurs de v\frflex{i}n me chansonnent.
\item Et moi, je te pr\frflex{i}e, Seigneur :\psstar{} c’est l’he\frflex{u}re de ta grâce.
\item Dans ton grand amour, Die\frflex{u}, réponds-moi,\psstar{} par ta vérit\frflex{é} sauve-moi.
\item Tire-moi de la boue, sin\frflex{o}n je m’enfonce :\psstar{} que j’échappe à ceux qui me haïssent, à l’ab\frflex{î}me des eaux.
\item Que les flots ne me submergent pas, que le go\frflex{u}ffre ne m’avale,\psstar{} que la gueule du puits ne se ferme p\frflex{a}s sur moi.